\scanpage{5}
\mspage{4}
Exemple : Si $G$ est un groupe algébrique sur $k$, la catégorie $\mathrm{Modf}(G)$ des vectoriels sur lesquels $K$ opère. Item avec groupes proalgébriques (p. ex. profinis).

Cet exemple \deleted{est} essentiellement universel, comme on va voir.

$\otimes$-foncteurs. Foncteur\unsure{s} fibre\unsure{s} \inserted{sur $K$, -- sur $K'$, \ldots.} Considérons un couple $(C,F)$ \deleted{\unsure{$s/F$|formé}} d'une $\otimes$-catégorie sur $K$, avec foncteur fibre \unsure{ntr.|etc.|rat.} /$K$. À cette situation \deleted{\unsure{est}} on peut associer un groupe \deleted{alg} affine $G$ sur $K$ (ou pro-groupe pro-affine sur $K$), et une équivalence $(C,F)\simeq(\mathrm{Modf}(G),\text{oubli})$.
% slanted
\margin{$G(A)=\operatorname{Aut}F_A$\quad $F_A(M)=F(M)\otimes_K A$ \deleted{\unsure{considéré}} comme foncteur fibre à valeurs dans $\mathrm{Mod}(A)$}
De façon précise, associant à chaque $G$ le couple $\bigl((\mathrm{Mod}(\unsure{f|G}),G),\text{oubli}\bigr)$, on trouve un foncteur
\diagram{\begin{tikzcd}[ampersand replacement=\&,column sep=large]
\text{groupe affine sur }K \arrow[r] \& \begin{gathered}\otimes\ \text{catégorie avec foncteur}\\\text{fibre donné, avec }\otimes\\\text{-foncteurs mod isom,}\\\text{compatibles avec }F\end{gathered}
\end{tikzcd}}
et ce foncteur est pleinement fidèle.
\[
\begin{aligned}
C\ \text{``de type fini''}&\iff G\text{ de type fini}\\
C\ \tfrac12\text{-simple}&\underset{\inserted{\text{si }K\text{ de car. nulle}}}{\iff}\deleted{\unsure{G\ \text{fini réductif}}}\ G\text{ proréductif}.\\
&\phantom{\iff}\deleted{\unsure{\text{fibre algébrique}}}
\end{aligned}
\]
\uline{NB} Possibilité d'une définition des autres op. tensorielles telles $\wedge^{\unsure{i}}$, $\operatorname{Sym}^i$ (au fait, on voit que c'est indépendant de l'existence et du choix de $F$)

Si on a un foncteur fibre, tout autre foncteur fibre $F'$ s'obtient en ``tordant'' avec un pro-torseur $P$ sous $G$, $P=(P_i)$ :
\[
F'=P\times^G F
\]
% slanted
\margin{\uuline{à prouver} \unsure{Revient|Reviens|Remarque} : affirmer que si $C$ de type fini, deux foncteurs fibres deviennent isom. sur ext. finie}
% archive sheet number: 4

